% Source: Wald GR, physical PDF pp. 250--254
%         (printed pp. 237--241).
% Coverage: Section 9.5, Singularity Theorems; no numbered equations.
% Figures/tables/footnotes: Figure 9.6; four footnotes.
% Status: source-reviewed and compile-clean.
% Uncertainties: none.

\subsection{Singularity Theorems}\label{sec:9.5}

\setcounter{theorem}{0}

We now have developed the machinery required to prove some of the singularity
theorems. We shall give a complete proof of three theorems which establish the
existence of singularities in the sense of timelike or null geodesic incompleteness
under conditions relevant to cosmology and gravitational collapse. Then we shall
simply quote a fourth theorem that establishes the existence of singularities under
significantly weaker hypotheses.

The first theorem we shall prove can be interpreted as showing that if the universe is
globally hyperbolic and at one instant of time is expanding everywhere at a rate
bounded away from zero, then the universe must have begun in a singular state a
finite time ago.

\begin{theorem}\label{thm:9.5.1}
Let \((M,g_{ab})\) be a globally hyperbolic spacetime with
\(R_{ab}\xi^a\xi^b\geq0\) for all timelike \(\xi^a\), which will be the case if
Einstein's equation is satisfied with the strong energy condition holding for matter.
Suppose there exists a smooth (or, at least \(C^2\)) spacelike Cauchy surface
\(\Sigma\) for which the trace of the extrinsic curvature (for the past directed
normal geodesic congruence) satisfies \(K\leq C<0\) everywhere, where \(C\) is a
constant. Then no past directed timelike curve from \(\Sigma\) can have length
greater than \(3/\lvert C\rvert\). In particular, all past directed timelike geodesics
are incomplete.
\end{theorem}

\begin{proof}
Suppose there were a past directed timelike curve, \(\lambda\), from \(\Sigma\) with
length greater than \(3/\lvert C\rvert\). Let \(p\) be a point on \(\lambda\) lying
beyond length \(3/\lvert C\rvert\) from \(\Sigma\). By Theorem~9.4.5, there exists a
maximum length curve \(\gamma\) from \(p\) to \(\Sigma\), which, clearly, also must
have length greater than \(3/\lvert C\rvert\). By Theorem~9.4.3, \(\gamma\) must be a
geodesic with no conjugate point between \(\Sigma\) and \(p\). However, this
contradicts Proposition~9.3.4 which states that \(\gamma\) must have a conjugate
point between \(\Sigma\) and \(p\). Therefore, the original curve \(\lambda\) cannot
exist.
\end{proof}

The strongest unwanted hypothesis in Theorem~9.5.1 is that the universe be globally
hyperbolic. Indeed, given only Theorem~9.5.1, it might seem more reasonable to
conclude that an everywhere expanding universe must fail to be globally hyperbolic
rather than that it must be singular. Since Theorem~9.4.5 plays a critical role in the
proof, it might appear that the assumption of global hyperbolicity cannot be
eliminated. However, we now shall prove a theorem due to Hawking (1967) which does
remove this assumption. The main price paid for this removal is the additional
hypothesis that \(\Sigma\) be compact (i.e., the universe is ``closed'') and the
% SOURCE ERRATUM: printed ``past direct''; corrected to ``past directed.''
significantly weakened conclusion that only at least one past directed timelike
geodesic (rather than all past directed timelike curves) must be incomplete.

\begin{theorem}\label{thm:9.5.2}
Let \((M,g_{ab})\) be a strongly causal\footnote{It has been shown by Hawking
(1967) that the assumption of strong causality can be eliminated (see problem~3).}
spacetime with \(R_{ab}\xi^a\xi^b\geq0\) for all timelike \(\xi^a\), as will be the
case if Einstein's equation is satisfied with the strong energy condition holding for
matter. Suppose there exists a compact, edgeless, achronal, smooth spacelike
hypersurface \(S\) such that for the past directed normal geodesic congruence from
\(S\) we have \(K<0\) everywhere on \(S\). Let \(C\) denote the maximum value of
\(K\), so \(K\leq C<0\) everywhere on \(S\). Then at least one inextendible past
directed timelike geodesic from \(S\) has length no greater than \(3/\lvert C\rvert\).
\end{theorem}

\begin{proof}
Suppose all past directed inextendible timelike geodesics from \(S\) had length
greater than \(3/\lvert C\rvert\). Since the spacetime
\((\operatorname{int}[D(S)],g_{ab})\) satisfies the hypotheses of Theorem~9.5.1, all
inextendible past directed timelike geodesics from \(S\) must leave
\(\operatorname{int}[D(S)]\). Since \(H(S)\) is the boundary of \(D(S)\) (see
Proposition~8.3.6), all such geodesics must intersect \(H^-(S)\) before their length
becomes greater than \(3/\lvert C\rvert\). In particular, this implies that
\(H^-(S)\ne\varnothing\). We shall prove that \(H^-(S)\) must be compact and then
show that this leads to a contradiction.

The key step in proving compactness of \(H^-(S)\) is the demonstration that for each
\(p\in H^-(S)\) there exists a maximum length orthogonal geodesic from \(S\) to
\(p\). First, the length of any causal curve from \(S\) to \(p\in H^-(S)\) is bounded
from above by \(3/\lvert C\rvert\), so the least upper bound, \(\tau_0\), of the length
of all causal curves from \(S\) to \(p\) exists. We wish to find an orthogonal
geodesic from \(S\) to \(p\) with length \(\tau_0\). Let \(\{\lambda_n\}\) be a
sequence of timelike curves from \(S\) to \(p\) such that
\[
  \lim_{n\to\infty}\tau[\lambda_n]=\tau_0 .
\]
Choose \(q_n\in\lambda_n\) with \(q_n\ne p\) such that the sequence
\(\{q_n\}\) converges to \(p\) (see Figure~\ref{fig:9.6}). Since \(q_n\in I^+(p)\),
we have \(q_n\in\operatorname{int}[D^-(S)]\). Hence, by Theorem~9.4.5, there exists
a normal geodesic \(\gamma_n\) from \(S\) to \(q_n\) which maximizes the length of all
causal curves from \(S\) to \(q_n\). Clearly, we have
\[
  \lim_{n\to\infty}\tau[\gamma_n]=\tau_0 .
\]
Let \(r_n\) be the intersection point of \(\gamma_n\) with \(S\). Since \(S\) is
compact, there exists an accumulation point \(r\) of the sequence \(\{r_n\}\). Let
\(\gamma\) be the geodesic normal to \(S\) originating from \(r\). Then, because of
the continuous dependence of geodesics on their initial point and tangent vector,
\(\gamma\) must intersect \(H^-(S)\) at \(p\) and
\[
  \tau[\gamma]=\lim_{n\to\infty}\tau[\gamma_n]=\tau_0 .
\]
Thus, we have found the desired timelike geodesic orthogonal to \(S\) which
maximizes the length from \(S\) to \(p\).

To prove the compactness of \(H^-(S)\), we show that every sequence \(\{p_n\}\) in
\(H^-(S)\) has an accumulation point \(p\in H^-(S)\). Let
\(\{\widetilde\gamma_n\}\) be a sequence of maximum length orthogonal geodesics from
\(S\) to \(p_n\).

% Source: Wald GR, physical PDF p. 252 (printed p. 239), Figure 9.6.
% Coverage: Spacetime diagram used in the proof of Theorem 9.5.2.
% Figures/tables/footnotes: Figure 9.6.
% Status: source-reviewed and compile-clean.
% Uncertainties: none.

\begingroup
\renewcommand{\thefigure}{9.6}
\begin{figure}[tb]
  \centering
  \begin{tikzpicture}[x=1cm,y=1cm,>=Stealth,thick,line cap=round]
    \draw (-2.72,4.78) .. controls (-1.75,5.42) and (-.72,4.88) .. (.52,4.72)
      .. controls (1.66,4.58) and (2.38,4.82) .. (3.15,5.00);
    \node[right] at (2.02,5.08) {\(S\)};
    \draw (-2.42,1.40) -- (0,-.08) -- (2.42,1.40);
    \node[below] at (0,-.16) {\(H^-(S)\)};

    % The timelike curves lambda_n all run from S to the same p on H^-(S).
    % Each marked q_n lies on its corresponding lambda_n and the sequence
    % approaches p, as required in the proof of Theorem 9.5.2.
    \coordinate (p) at (-1.16,.63);
    \coordinate (qone) at (-1.62,1.50);
    \coordinate (qtwo) at (-1.32,1.17);
    \coordinate (qthree) at (-1.11,.91);
    \draw (p) .. controls (-1.35,.90) and (-1.51,1.25) .. (qone)
      .. controls (-1.82,2.22) and (-1.78,3.75) .. (-1.56,5.02);
    \draw (p) .. controls (-1.24,.82) and (-1.30,1.03) .. (qtwo)
      .. controls (-1.51,2.26) and (-1.33,3.77) .. (-.96,4.92);
    \draw (p) .. controls (-1.15,.74) and (-1.13,.84) .. (qthree)
      .. controls (-1.08,2.22) and (-.79,3.65) .. (-.34,4.82);

    % For every q_n, the maximizing normal geodesic gamma_n joins q_n to
    % its own intersection r_n with S.
    \coordinate (rone) at (-1.24,4.98);
    \coordinate (rtwo) at (-.68,4.88);
    \coordinate (rthree) at (-.10,4.78);
    \draw[dashed] (rone) .. controls (-1.30,3.57) and (-1.35,2.24) .. (qone);
    \draw[dashed] (rtwo) .. controls (-.79,3.51) and (-.98,2.11) .. (qtwo);
    \draw[dashed] (rthree) .. controls (-.26,3.42) and (-.61,1.91) .. (qthree);

    \fill (p) circle (2pt) node[below left] {\(p\)};
    \foreach \q in {qone,qtwo,qthree} {\fill (\q) circle (1.55pt);}
    \foreach \r in {rone,rtwo,rthree} {\fill (\r) circle (1.55pt);}
    \node[right] at (-.45,1.38) {\(q_n\)};
    \draw[->] (-.52,1.34) -- (qtwo);
    \node[above] at (-.60,5.36) {\(r_n\)};
    \draw[->] (-.72,5.27) -- (rtwo);
    \draw[->] (-.50,5.27) -- (rthree);
    \draw[->] (-2.38,3.32) to[out=5,in=175] (-1.64,3.49);
    \node[left] at (-2.34,3.12) {\(\lambda_n\)};
    \draw[->] (1.32,3.38) to[out=178,in=10] (-.31,3.30);
    \node[right] at (1.38,3.43) {\(\gamma_n\)};
    \node[right] at (.82,2.42) {\(D^-(S)\)};
    % Reserve trim clearance above the artwork when this float lands at page top.
    \path[draw=none] (0,6.4) circle[radius=.1pt];
  \end{tikzpicture}
  \caption{A spacetime diagram showing a construction used in the proof of
  theorem~9.5.2.}
  \label{fig:9.6}
\end{figure}
\endgroup


We, now, in essence, repeat the argument given at the end of the previous paragraph.
Let \(\widetilde r_n\) be the intersection point of \(\widetilde\gamma_n\) with \(S\),
and let \(\widetilde r\) be an accumulation point of \(\{\widetilde r_n\}\). Let
\(\widetilde\gamma\) be the geodesic starting from \(\widetilde r\) orthogonal to
\(S\), and let \(p\) be the intersection point of \(\widetilde\gamma\) with
\(H^-(S)\). Then \(p\) is an accumulation point of \(\{p_n\}\). Thus, \(H^-(S)\) is
compact.

However, since \(\operatorname{edge}(S)=\varnothing\), by Theorem~8.3.5,
\(H^-(S)\) contains a future inextendible null geodesic. Since \((M,g_{ab})\) is
strongly causal, by Lemma~8.2.1 this is impossible if \(H^-(S)\) is compact. Thus,
our assumption that all past directed inextendible timelike geodesics from \(S\) have
length greater than \(3/\lvert C\rvert\) has led to a contradiction.
\end{proof}

The previous two theorems established timelike geodesic incompleteness in cosmological
contexts. The next theorem proves null geodesic incompleteness in a context relevant
to gravitational collapse. A compact, two-dimensional, smooth spacelike
submanifold, \(T\), having the property that the expansion, \(\theta\), of \emph{both}
sets (i.e., ``ingoing'' and ``outgoing'') of future directed null geodesics
orthogonal to \(T\) is everywhere negative is called a \emph{trapped surface}. In the
extended Schwarzschild solution, all spheres inside the black hole (region II of
Figure~6.9) are trapped surfaces. As will be discussed further in chapter~12, it
follows from this fact that trapped surfaces must form in any gravitational collapse
whose initial conditions are sufficiently close to initial conditions for spherical
collapse. The next theorem---which, historically, was the first general singularity
theorem to be proven (Penrose 1965a)---shows that, under some further hypotheses, a
singularity must occur after a trapped surface has formed.

\begin{theorem}\label{thm:9.5.3}
Let \((M,g_{ab})\) be a connected, globally hyperbolic spacetime with a noncompact
Cauchy surface \(\Sigma\). Suppose \(R_{ab}k^ak^b\geq0\) for all null \(k^a\), as
will be the case if \((M,g_{ab})\) is a solution of Einstein's equation with matter
satisfying the weak or strong energy condition. Suppose, further, that \(M\) contains
a trapped surface \(T\). Let \(\theta_0<0\) denote the maximum value of \(\theta\)
for both sets of orthogonal geodesics on \(T\). Then at least one inextendible future
directed orthogonal null geodesic from \(T\) has affine length no greater than
\(2/\lvert\theta_0\rvert\).
\end{theorem}

\begin{proof}
Suppose all future directed null geodesics from \(T\) have affine length
\(\geq2/\lvert\theta_0\rvert\). Then we may define the map
\(f_+:T\times[0,2/\lvert\theta_0\rvert]\longrightarrow M\) by setting \(f(q,a)\)
equal to the point of \(M\) lying at affine parameter \(a\) on the ``outgoing'' null
geodesic normal to \(T\) starting at \(q\). Similarly, we define \(f_-\) for the
``ingoing'' geodesics. Since \(T\times[0,2/\lvert\theta_0\rvert]\) is a compact set
and \(f_+\) and \(f_-\) are continuous, the images of \(f_+\) and \(f_-\) and hence
their union
\[
  A=f_+\bigl\{T\times[0,2/\lvert\theta_0\rvert]\bigr\}
    \cup f_-\bigl\{T\times[0,2/\lvert\theta_0\rvert]\bigr\}
\]
must also be compact (see Appendix~A). However, by Proposition~9.3.9 and
Theorem~9.3.11, \(\partial I^+(T)\) is a subset of \(A\), and since
\(\partial I^+(T)\) is closed, we conclude that \(\partial I^+(T)\) is compact.

We show, now, that compactness of \(\partial I^+(T)\) contradicts the existence of a
noncompact Cauchy surface \(\Sigma\). Using Lemma~8.1.1, we choose a smooth timelike
vector field \(t^a\) on \(M\). Since \(\partial I^+(T)\) is achronal, each integral
curve of \(t^a\) can intersect \(\partial I^+(T)\) at most once, while every integral
curve of \(t^a\) intersects \(\Sigma\) precisely once. Thus, we may define a map
\(\psi:\partial I^+(T)\longrightarrow\Sigma\) by following the integral curves of
\(t^a\) from \(\partial I^+(T)\) to \(\Sigma\). Let \(S\subseteq\Sigma\) denote the
image, \(\psi[\partial I^+(T)]\), of \(\partial I^+(T)\) under \(\psi\), and let \(S\)
be given the topology induced by \(\Sigma\). Then
\(\psi:\partial I^+(T)\longrightarrow S\) is a homeomorphism. Since
\(\partial I^+(T)\) is compact, so is \(S\), and, hence, viewed as a subset of
\(\Sigma\), \(S\) must be closed. On the other hand, since \(\partial I^+(T)\) is a
\(C^0\)-manifold (see Theorem~8.1.3), each point of \(\partial I^+(T)\) has a
neighborhood homeomorphic to an open ball in \(\mathbb R^3\). Since \(\psi\) is a
homeomorphism, the same property holds for \(S\), and, hence, viewed as a subset of
\(\Sigma\), \(S\) must be open. However, since \(M\) is connected, \(\Sigma\) also
must be connected (see Theorem~8.3.14). Therefore, since
\(\partial I^+(T)\ne\varnothing\), we must have \(S=\Sigma\). However, this is
impossible since \(S\) is compact but \(\Sigma\) is noncompact.
\end{proof}

Again, Theorem~9.5.3 contains the unwanted hypothesis that \((M,g_{ab})\) is
globally hyperbolic. However, with some additional assumptions, this hypothesis
again can be eliminated by arguments similar to those given in the proof of
Theorem~9.5.2. We shall not attempt to give these arguments here, but will merely
quote a singularity theorem of Hawking and Penrose (1970) which eliminates most of
the unwanted assumptions, referring the reader to that reference or Hawking and
Ellis (1973) for a proof.

\begin{theorem}\label{thm:9.5.4}
Suppose a spacetime \((M,g_{ab})\) satisfies the following four conditions.
\begin{enumerate}
  \item \(R_{ab}v^av^b\geq0\) for all timelike and null \(v^a\), as will be the case
  if Einstein's equation is satisfied with the strong energy condition holding for
  matter.
  \item The timelike and null generic conditions are satisfied (see Section~9.3).
  \item No closed timelike curve exists.
  \item At least one of the following three properties holds:
  \begin{enumerate}
    \item \((M,g_{ab})\) possesses a compact achronal set without edge [i.e.,
    \((M,g_{ab})\) is a closed universe].
    \item \((M,g_{ab})\) possesses a trapped surface, or
    \item there exists a point \(p\in M\) such that the expansion of the future (or
    past) directed null geodesics emanating from \(p\) becomes negative along each
    geodesic in this congruence.
  \end{enumerate}
\end{enumerate}
Then \((M,g_{ab})\) must contain at least one incomplete timelike or null geodesic.
\end{theorem}

As compared with Theorem~9.5.2, the above theorem adds only the hypothesis that the
generic conditions are satisfied. However, it entirely eliminates the assumption that
the universe is expanding everywhere. Similarly, as compared with Theorem~9.5.3, the
above theorem adds only the generic condition plus the condition that
\(R_{ab}\xi^a\xi^b\geq0\) for all timelike vectors. However, it eliminates entirely
the assumption that \((M,g_{ab})\) be globally hyperbolic. Thus, Theorem~9.5.4 has
much wider applicability than the three previous theorems proven above. On the other
hand, the conclusions of Theorem~9.5.4 are slightly weaker in that no information is
provided concerning which timelike or null geodesic is incomplete.

Gannon (1975) has strengthened Theorem~9.5.4 by adding a fourth alternative to
condition 4, namely that \((M,g_{ab})\) possesses a closed, achronal, edgeless set
\(S\) which is non-simply connected\footnote{See chapter~13 for the definition of
simply connected.} and is ``regular near infinity.'' Here ``regular near infinity''
means that \(S\) can be expressed as a union of a nested family of sets, \(W_i\)
(i.e., each \(W_i\) satisfies \(W_i\subseteq W_{i+1}\)), such that (i) each \(W_i\) is
compact, (ii) its boundary, \(\partial W_i\), in \(S\) is homeomorphic to a 2-sphere,
(iii) \(S-\operatorname{int}(W_i)\) is homeomorphic to
\(S^2\times\mathbb R^+\), where \(\mathbb R^+=[0,\infty)\), and (iv) the expansion,
\(\theta\), of the inward directed null geodesics orthogonal to each \(W_i\) is
negative everywhere on \(W_i\). Except for the compactness\footnote{A singularity
theorem applicable to the case where the \(W_i\) are noncompact is given by Gannon
(1976).} of the \(W_i\), the four conditions defining asymptotic regularity are
satisfied by all asymptotically flat hypersurfaces in asymptotically flat spacetimes
(see chapter~11). Thus, in essence, Gannon's theorem shows that an asymptotically
flat spacetime satisfying conditions (1)--(3) of Theorem~9.5.4 and which
``initially'' has appropriately nontrivial topology must develop a singularity. We
refer the reader to Gannon (1975, 1976) for further details and results.

Theorem~9.5.4 gives us strong reason to believe that our universe is singular. As
discussed in chapter~5, the observational evidence strongly suggests that our
universe---or, at least, the portion of our universe within our causal past---is well
described by a Robertson--Walker model (differing not too greatly from a \(k=0\)
model), at least back as far as the decoupling time of matter and radiation. However,
in these models, the expansion of the past directed null geodesics emanating from the
event representing us at the present time becomes negative at a much more recent time
than the decoupling time. Thus, there is strong reason to believe that condition 4(c)
of Theorem~9.5.4 is satisfied in our universe.\footnote{This conclusion also can be
drawn from much weaker assumptions; see chapter~10 of Hawking and Ellis (1973).}
Since we expect that conditions (1)--(3) also are satisfied, it appears that our
universe must be singular. Thus, it appears that we must confront the breakdown of
classical general relativity expected to occur near singularities if we are to
understand the origin of our universe.
